% ---
% §1.3 Justificación. Reemplaza las líneas 607–631 de Docs/Technical_Report.tex.
% Cambios respecto a TT1 (auditoría T3, §1.3):
%   - \cite de Azevedo et al. (estaba nombrado sin cita) y de Lee et al. / Clark et al. para el CBIS-DDSM;
%   - «circuitos de hasta 16 qubits» → capacidad medida en este trabajo (H-051);
%   - fuera «si demuestran mejorar … clasificadores más generalizables» (presupone el resultado)
%     y «que la literatura menciona pero no sistematiza» (sin fuente);
%   - la relevancia social remite a §1.1 en lugar de repetir cifras;
%   - «datasets predominantes en la literatura» → remite a la tabla comparativa de §2.
% Decisiones y hallazgos que lo respaldan: D-012, D-014, D-019, D-032, D-033; H-051.
% ---

\subsection{Justificación}
\label{sec:justificacion}

\textbf{Relevancia científica}

% respaldo: «Furthermore, results on small datasets cannot be used to reason about larger datasets» (Bowles2024Benchmarking, p. 21)
% respaldo: «There is hence a need to study how results scale to larger datasets, which typically pushes the number of qubits to the limits of what is possible today.» (Bowles2024Benchmarking, p. 22)
El aprendizaje automático cuántico necesita evidencia empírica sobre datos reales. Una parte
importante de su evaluación se ha hecho sobre problemas de prueba construidos para ilustrar
propiedades teóricas \cite{AlvarezEstevez2024QuantumKernelBench}, y los resultados obtenidos con
conjuntos de datos pequeños no permiten razonar sobre conjuntos mayores \cite{Bowles2024Benchmarking}.
Llevar estas técnicas a un dominio de imagen médica, con miles de lesiones y la variabilidad propia de
las mamografías digitalizadas, es la forma de saber cómo se comportan fuera de esas condiciones.

% respaldo: «As future work, an interesting approach could be to separate the dataset into masses and calcifications and only then perform classification.» (Azevedo2022QuantumTransfer, p. 12)
% respaldo: «Lastly, the best way to know how well this model performs would be to see how the data appears in our Hilbert space.» (Azevedo2022QuantumTransfer, p. 12)
El presente trabajo contribuye a esa evidencia de tres formas. Primero, documenta el comportamiento
de \textit{ZZFeatureMap} y \textit{PauliFeatureMap} de Qiskit aplicados a características radiómicas
extraídas de mamografías reales en formato DICOM. Segundo, evalúa las representaciones por su
separabilidad de clases, y no solo por el desempeño del clasificador final, con métricas definidas
sobre los kernels y sobre los propios vectores de representación: el alineamiento kernel-objetivo, la
diferencia geométrica, el índice de Davies-Bouldin y la razón de Fisher (sección~\ref{sec:metricas}).
Tercero, retoma dos de las líneas de trabajo futuro que dejan Azevedo et al.
\cite{Azevedo2022QuantumTransfer}: analizar por separado masas y calcificaciones, y observar cómo se
distribuyen los datos en el espacio de Hilbert. Así, los resultados se sitúan en una línea de
investigación acumulativa en aprendizaje automático cuántico para el diagnóstico médico.

\textbf{Relevancia social}

El cáncer de mama es la principal causa de muerte por tumores malignos en las mujeres de México, y su
detección temprana determina en gran medida el desenlace clínico (sección~\ref{sec:introduccion}).
Las herramientas de prediagnóstico asistido pueden ampliar la capacidad de un sistema de salud en el
que la disponibilidad de mastógrafos y de especialistas en radiología mamaria no cubre la demanda. Para
orientar ese desarrollo es necesario saber qué técnicas aportan y cuáles no: si las representaciones
cuánticas no mejoran la separabilidad de las lesiones, ese resultado, medido con rigor, también indica
dónde no conviene invertir el esfuerzo de investigación. Aunque el presente trabajo es un prototipo
académico y no un dispositivo clínico, sus resultados se suman al conocimiento que eventualmente
informará el desarrollo de herramientas diagnósticas para sistemas de salud con recursos limitados,
como el mexicano.

\textbf{Viabilidad técnica}

% respaldo: «The data were split into a training set and a testing set based on the BI-RADS category.» (Lee2017CBISDDSM, p. 5)
% fuente: D-013, H-032, H-051; B1_benchmark_arquitectura.csv, B1_benchmark_kernel_proyeccion.csv, 8_oe6_qubits.csv, 8_oe6_resumen.csv
% VERIFICAR: licencia CC BY 3.0 del CBIS-DDSM en TCIA (la ficha de Lee no la respalda con cita literal).
La viabilidad del trabajo descansa en tres recursos. El primero es el dataset CBIS-DDSM, una versión
curada del DDSM con diagnóstico confirmado por patología, disponible públicamente en The Cancer
Imaging Archive \cite{Lee2017CBISDDSM, Clark2013TCIA}. En la versión que distribuye TCIA, reúne 3\,568 lesiones de 1\,566 pacientes, con
una partición oficial de entrenamiento y prueba estratificada por categoría BI-RADS y metadatos
clínicos estructurados, lo que evita procesos propios de adquisición, anonimización o validación
clínica. El segundo es la simulación clásica de los circuitos en el equipo de desarrollo, un
procesador Apple M3 con 16 GiB de memoria. El número de qubits se fijó en 12 tras un estudio de
costos: con la implementación inicial, que evaluaba cada valor del kernel con un circuito propio, la
etapa cuántica completa pasaba de unas 3~h con 12 qubits a unas 41~h con 16, y el diseño requiere
decenas de kernels. Después se adoptó una implementación que calcula el vector de estado de cada
lesión una sola vez y obtiene el kernel de sus productos internos; con ella se comprobó que el límite
físico de la máquina está más arriba, pues se simularon estados de hasta 28 qubits
(sección~\ref{sec:resultados}). Estas cifras corresponden a la simulación ideal del vector de estado,
sin ruido, que además permite aislar el efecto de la transformación cuántica sin la interferencia del
ruido de los dispositivos físicos. El tercero es un ecosistema de
programación de código abierto, documentado y mantenido (Qiskit, PyTorch, PyRadiomics, scikit-learn),
que permite a otros investigadores reproducir el trabajo.

\textbf{Novedad y posicionamiento frente a trabajos previos}

La aportación del trabajo se define por tres diferenciadores. El primero es el dominio de datos:
opera sobre mamografías DICOM del CBIS-DDSM y no sobre datos clínicos tabulares, como Wisconsin, SEER o
GBSG, empleados en otros trabajos de aprendizaje automático cuántico en oncología
(cuadro~\ref{tab:estado_arte_v2}). El segundo es el papel del componente cuántico: el circuito se usa
exclusivamente como transformación de la representación, con parámetros fijos, y no como una capa que
se entrena dentro de una red, lo que permite analizar de forma aislada su efecto sobre la
separabilidad de las clases. El tercero es el análisis de viabilidad: las restricciones prácticas del
simulador se convierten en un objetivo de investigación explícito y medible, con datos cuantitativos
sobre el tiempo de ejecución, el número de qubits simulables, el número de mediciones que requeriría
un procesador real y el costo de entrenar el circuito. En conjunto, estos diferenciadores sitúan el
trabajo como una contribución experimental a una línea de investigación activa y como una referencia
metodológica para aplicar codificaciones cuánticas a la imagen médica.
